\chapter{The financial problem}\label{cap:problema_financiero}

This chapter sets out, in the language of physics and without assuming
stochastic calculus, the problem this work solves: which contracts are
priced, why their price is determined, and why path-dependent options
call for a different method. For the proofs we refer to
\cite{tfgm, shreve, hull}.

\section{Contracts and payoffs}\label{sec:contratos}

\subsection{European options}\label{subsec:europeas}

We consider a market with two assets. The first is a risky
asset---a stock, for example---whose price $S_t$ fluctuates
randomly; in the following chapters it will play the role of the
coordinate of a particle. The second is a riskless bank account, whose
balance grows deterministically, $B_t = B_0\, e^{rt}$, with a constant
interest rate $r$. One euro received at $T$ is therefore worth
$e^{-r\tau}$ euros at time $t$, with $\tau = T - t$: the discount factor
converts amounts from one date to another, like a change of units, and
it only makes sense to compare prices referred to the same date.

A \emph{derivative} is a contract whose payoff depends on the price of
another asset, called the \emph{underlying}. The simplest is the
\emph{European call option} with exercise price (\emph{strike}) $K$ and
maturity $T$: it gives its holder the right, but not the obligation, to
buy one unit of the underlying at time $T$ at the price $K$. If
$S_T > K$, the holder exercises and gains $S_T - K$; if $S_T \leq K$,
the holder does not exercise and gains nothing. The \emph{payoff} at
$T$ is
\begin{equation}\label{eq:pago_call}
  g(S_T) \;=\; (S_T - K)^+ \;\eqdef\; \max(S_T - K,\, 0).
\end{equation}
The European put option pays $(K - S_T)^+$, and in general we call
European any option that pays a function $g(S_T)$ of the price at
maturity. Its structural feature is that the payoff depends only on the
final state: in the language of physics, $g$ is an observable measured
at $t = T$ on the position $S_T$, and the history that led there is
irrelevant. To price it, it suffices to know the law of a single random
variable, $S_T$.

The problem is to determine how much the right to receive $g(S_T)$ at
$T$ is worth today, at $t < T$. It is worth explaining in advance why a
model of the underlying is needed. If the payoff were linear, $S_T - K$
(a \emph{forward} contract), its price would be fixed without any
assumption on the dynamics of $S$: buying one unit of the underlying
today and borrowing $K e^{-r\tau}$ reproduces that payoff in every
scenario, so the contract is worth $S_t - K e^{-r\tau}$. The call, by
contrast, has a kink at $S_T = K$: it is convex, and its value depends
on the whole distribution of $S_T$ and not only on its mean. That is why
the volatility of the underlying enters, and why a model is needed
(Section~\ref{sec:modelo_bs}).

The options studied in this work, in contrast, pay a \emph{path
functional}: the payoff at $T$ is $\Phi[S(\cdot)]$, with
$S(\cdot) = \{S_t\}_{0 \leq t \leq T}$, and the European option is the
particular case $\Phi[S(\cdot)] = g(S_T)$. The two families in this work
are Asian options, which depend on an integral along the path, and
barrier options, which depend on its extreme values.

\subsection{Asian options}
\label{subsec:asian_concepto}\label{sec:asian_definicion}
\label{subsec:asian_promedios}

An Asian option is a derivative whose payoff depends on the
\emph{average} value of the underlying over a monitoring period, instead
of its terminal value: the average is harder to manipulate than a single
observation and, since it has less variance, the option is cheaper than
the equivalent European one. Over the period $[0, T]$, the two canonical
averages are the arithmetic and the geometric one:
\begin{align}
  A_T &\;\eqdef\; \frac{1}{T}\int_0^T S_t\, dt,
  \label{eq:promedio_aritmetico}\\
  G_T &\;\eqdef\; \exp\!\left(\frac{1}{T}\int_0^T \log S_t\, dt\right).
  \label{eq:promedio_geometrico}
\end{align}
The canonical case, and the one priced in this work, is the
\emph{fixed-strike call}, which pays $(A_T - K)^+$ or $(G_T - K)^+$
depending on the average; the put and the floating-strike variants,
which compare $S_T$ with the average, are treated in the same way. In
practice the average is taken over a finite number of dates; here
continuous monitoring, which is its limit, is always used.

\medskip

\noindent\textbf{Linearity in the log-price.}
The observation that structures Chapters~\ref{cap:fuente_pared}
and~\ref{cap:espectro}\enprep{} appears on passing to the log-price $X_t = \log S_t$.
The logarithm of the geometric average is linear in the path; that of
the arithmetic average is not:
\begin{align}
  \log G_T &\;=\; \frac{1}{T}\int_0^T X_t\, dt,
  \label{eq:log_geom_lineal}\\
  \log A_T &\;=\; \log\!\left(\frac{1}{T}\int_0^T e^{X_t}\, dt\right).
  \label{eq:log_arit_no_lineal}
\end{align}
In the Black--Scholes model, $X_t$ is a Gaussian process
(Section~\ref{sec:modelo_bs}). A linear
combination of Gaussian variables is Gaussian, so $G_T$ is lognormal and
its option admits an elementary treatment; $A_T$, an integral of
exponentials of a Gaussian process, is not lognormal and has no
elementary closed-form density, and it is this nonlinearity that makes
the arithmetic case fundamentally harder. By the inequality between the
arithmetic and geometric means, a consequence of the convexity of the
exponential, $G_T \leq A_T$ on every path; for realistic volatilities
and maturities the difference is appreciable, and the arithmetic case
does not reduce to the geometric one.

\subsection{Barrier options}
\label{subsec:barrera_unica_payoff}\label{sec:barrera_definicion}
\label{subsec:doble_barrera_otros_tipos}

A barrier option pays like a European one, but only if the underlying
has, or has not, crossed a level $B$ fixed in the contract. The option
is deactivated (\emph{knock-out}) or activated (\emph{knock-in}) on
touching the barrier, which can lie below (\emph{down}) or above
(\emph{up}) the initial price. The canonical case is the
\emph{down-and-out call}, with $B < S_0$:
\begin{equation}\label{eq:payoff_DO_call}
  C_{DO}(T) \;=\; (S_T - K)^+ \cdot
  \1\!\left\{\min_{0 \leq t \leq T} S_t > B\right\}.
\end{equation}
The indicator cancels the payoff if the underlying touched $B$ at some
instant, whatever $S_T$ is: the path enters through its minimum. Since
every path either touches the barrier or does not,
$\1\{\min_t S_t > B\} + \1\{\min_t S_t \leq B\} = 1$, and the payoffs of
the \emph{down-and-out} and the \emph{down-and-in} with the same barrier
add up to that of the European option. The same holds for their prices,
because non-additive prices would allow a riskless profit
(\emph{in--out parity}):
\begin{equation}\label{eq:paridad_in_out}
  C_{DO}(t, S) + C_{DI}(t, S) \;=\; C_{\text{vanilla}}(t, S).
\end{equation}
It is therefore enough to price the \emph{out} options; the European one
is obtained in Chapter~\ref{cap:bs_operadores}\enprep. In the
\emph{double barrier} option two levels $B_d < B_u$ are fixed, and the
payoff is received only if the underlying stays in $(B_d, B_u)$
throughout the life of the contract: the \emph{double knock-out call}
pays $(S_T - K)^+\, \1\{B_d < \min_t S_t,\ \max_t S_t < B_u\}$.
Throughout this work the barrier has a constant level, is monitored
continuously, and no compensation (\emph{rebate}) is paid when the
option is knocked out.

\section{No arbitrage and replication in a one-period model}\label{sec:no_arbitraje}

The principle that fixes prices is that of \emph{no arbitrage}: there
is no portfolio that costs zero or less today, loses in no scenario and
gains in some scenario with positive probability. It is an equilibrium
condition: if such a portfolio existed, everyone would exploit it and
prices would adjust until it disappeared. Let us see its effect in the
simplest possible model \cite[Section~2.2]{tfgm}.

\medskip

\noindent\textbf{The model.}
There is a single period, of length $\Delta t$. The underlying moves
from $S_0$ to $S_0 u$ with probability $p$ or to $S_0 d$ with
probability $1 - p$, with $0 < d < u$, and the bank account multiplies
its balance by $e^{r\Delta t}$. No arbitrage requires
\begin{equation}\label{eq:no_arbitraje_un_periodo}
  d \;<\; e^{r\Delta t} \;<\; u .
\end{equation}
If $e^{r\Delta t} \geq u$, selling the asset short (borrowing it,
selling it, and returning it at the end) and depositing the proceeds never loses
and gains if the asset falls; if $e^{r\Delta t} \leq d$, borrowing to
buy the asset never loses and gains if it rises.

\medskip

\noindent\textbf{Replication.}
An option that expires at the end of the period pays $V_u$ if the asset
rises and $V_d$ if it falls; for the call, $V_u = (S_0 u - K)^+$ and
$V_d = (S_0 d - K)^+$. We form a portfolio with $\Delta$ units of the
underlying and $\vartheta$ euros in the account, and require it to pay
the same as the option in both states:
$\Delta\, S_0 u + \vartheta\, e^{r\Delta t} = V_u$ and
$\Delta\, S_0 d + \vartheta\, e^{r\Delta t} = V_d$. This linear system
has a unique solution,
\begin{equation}\label{eq:delta_un_periodo}
  \Delta \;=\; \frac{V_u - V_d}{S_0 u - S_0 d},
  \qquad
  \vartheta \;=\; e^{-r\Delta t}\,\frac{u V_d - d V_u}{u - d}.
\end{equation}
The portfolio costs $\Delta S_0 + \vartheta$ today. If the option traded
at a different price, buying the cheaper of the two and selling the
more expensive one would yield a riskless profit with no initial
investment, so no arbitrage fixes
\begin{equation}\label{eq:precio_un_periodo}
  V_0 \;=\; \Delta S_0 + \vartheta
  \;=\; e^{-r\Delta t}\bigl[\,q\, V_u + (1 - q)\, V_d\,\bigr],
  \qquad
  q \;\eqdef\; \frac{e^{r\Delta t} - d}{u - d}.
\end{equation}

\noindent\textbf{Interpretation.}
The probability $p$ does not appear in \eqref{eq:precio_un_periodo}.
The reason is geometric: with two possible states, the payoffs at the
end of the period form a two-dimensional vector space, and the
underlying and the account, with linearly independent payoffs
$(S_0 u, S_0 d)$ and $(e^{r\Delta t}, e^{r\Delta t})$, are a basis. Any
option is a vector in that space, and its price is fixed by linearity
when it is decomposed in the basis, without the likelihood of each
state playing any role. In particular, the price does not depend on how
much the underlying is expected to grow: in the continuous model, it
will not depend on its expected return $\mu$. Moreover, $\Delta$ in
\eqref{eq:delta_un_periodo} is a difference quotient,
$\Delta V / \Delta S$, the discrete version of the hedge
ratio $\Delta = \partial_S V$. With three states replication is no longer
possible in general (an \emph{incomplete} market,
\cite[Section~2.2.5]{tfgm}); the Black--Scholes model, the
$\Delta t \to 0$ limit of a model with many periods
\cite[Sections~2.3--2.4]{tfgm}, is complete.

\section{The risk-neutral probability}\label{sec:neutral_riesgo}

Condition \eqref{eq:no_arbitraje_un_periodo} is equivalent to
$0 < q < 1$, so $q$ can be read as a probability. We call $\Q$ the one
that assigns probability $q$ to the up move and $1 - q$ to the down
move. With it, \eqref{eq:precio_un_periodo} says that the price of the
option is its discounted payoff averaged with $\Q$,
\begin{equation}\label{eq:precio_q_un_periodo}
  V_0 \;=\; e^{-r\Delta t}\, \E_\Q[V_1],
\end{equation}
and a direct calculation shows what characterizes $\Q$:
\begin{equation}\label{eq:q_crecimiento_r}
  \E_\Q[S_1] \;=\; q\, S_0 u + (1 - q)\, S_0 d \;=\; S_0\, e^{r\Delta t}.
\end{equation}
Under $\Q$ the risky asset grows on average at the risk-free rate, like
the bank account: hence the name \emph{risk-neutral} probability. Put
differently, the discounted price $e^{-rt}S_t$ has the same expectation
at the end of the period as at the beginning. A process whose expected
future value, conditional on the present information, equals its
present value is a \emph{martingale}: a fair game, with no drift to
exploit. By \eqref{eq:precio_q_un_periodo}, the discounted price of the
option is also a martingale under $\Q$.

Two remarks prevent misunderstandings. The first: $\Q$ is not a
prediction of what the market will do. The real-world probabilities are in
$p$, which plays no role; $\Q$ is the only set of weights with which
averaging the payoffs is consistent with no arbitrage, and using these
weights is a computational device. The second: $\Q$ and the real-world probability
$\PP$ assign probability zero to the same events (here to none, because
$0 < p, q < 1$). They change the weights of the scenarios, not which
scenarios are possible; this is what it means for them to be
\emph{equivalent}.

The result generalizes to the \emph{fundamental theorem of asset
pricing} \cite[Section~5.4]{shreve}: a market admits no arbitrage if
and only if there is a probability $\Q$ equivalent to $\PP$ under which
the discounted prices of all assets are martingales, and this $\Q$ is
unique if and only if the market is complete. The price of a replicable
derivative is then its discounted payoff averaged with $\Q$; for the
binomial model, the explicit construction is in
\cite[Sections~2.2--2.3]{tfgm}.

\section{The Black--Scholes model}\label{sec:modelo_bs}

The Black--Scholes model describes the underlying with the simplest
assumptions compatible with a positive, random price: constant
volatility $\sigma$ and rate $r$, paths without jumps, and continuous,
costless trading. Its ingredient is Brownian motion, of which only the
minimum is needed here; its physics (the random walk, the heat equation
and the Langevin equation) is developed in
Chapter~\ref{cap:fundamentos}\enprep.

\medskip

\noindent\textbf{Brownian motion and the rule $dW^2 = dt$.}
Brownian motion $W_t$ starts at $W_0 = 0$ and has continuous paths and
independent increments over disjoint intervals, with
\begin{equation}\label{eq:incremento_browniano}
  W_{t + \Delta t} - W_t \;\sim\; N(0, \Delta t).
\end{equation}
A typical increment is of size $\sqrt{\Delta t}$, much larger than
$\Delta t$ when $\Delta t$ is small: the paths are continuous but
nowhere differentiable, and Taylor expansions must be carried to second
order in $\Delta W$. This second-order term behaves in a remarkable way.
Each $(\Delta W)^2$ has mean $\Delta t$ and variance $2\Delta t^2$;
summed over the $t/\Delta t$ steps of $[0, t]$, they give a variable
with mean $t$ and variance $2t\,\Delta t$, which tends to zero. In the
limit, the sum of the squared increments is no longer random and equals
$t$ on every path (this is the \emph{quadratic variation},
\cite[Section~3.2.3]{tfgm}). In an infinitesimal expansion this is
summarized in the rule
\begin{equation}\label{eq:regla_dW2}
  (dW)^2 = dt, \qquad dW\, dt = 0, \qquad (dt)^2 = 0,
\end{equation}
where the last two mean that those terms are of higher order than $dt$.
It is the only new rule this chapter needs.

\medskip

\noindent\textbf{The model.}
The price of the underlying follows the equation
\begin{equation}\label{eq:bs_eds_P}
  dS_t \;=\; \mu\, S_t\, dt + \sigma\, S_t\, dW_t,
\end{equation}
the limit of $S_{t+\Delta t} = S_t\,(1 + \mu\,\Delta t + \sigma\,\Delta W)$:
the relative return $dS/S$ has a drift $\mu$ and a noise of amplitude
$\sigma$, the \emph{volatility}. For the log-price $X_t = \log S_t$, the
second-order Taylor expansion, with $(dS/S)^2 = \sigma^2\,dt$ by
\eqref{eq:regla_dW2}, gives
\begin{equation}\label{eq:logprecio_P}
  dX_t \;=\; \frac{dS_t}{S_t} - \frac{1}{2}\left(\frac{dS_t}{S_t}\right)^{\!2}
  \;=\; \Bigl(\mu - \frac{\sigma^2}{2}\Bigr)\, dt + \sigma\, dW_t .
\end{equation}
This is the Langevin equation of a Brownian particle subject to a
constant force, with drift $\mu - \sigma^2/2$ and diffusion coefficient
$\sigma^2/2$. Its coefficients are constant, so
$X_t = X_0 + (\mu - \sigma^2/2)\,t + \sigma W_t$ is Gaussian and $S_t$
is lognormal. The term $-\sigma^2/2$ comes from the concavity of the
logarithm: rising by 10\% and then falling by 10\% leaves the price at
$0.99$ times its initial value. Volatility thus slows down the typical
growth of a path, measured by the median $S_0\,e^{(\mu - \sigma^2/2)t}$,
even though the mean is still $S_0\,e^{\mu t}$. The same expansion,
applied to a function $f(t, S)$, gives \emph{Itô's rule}
\cite[Section~4.3]{tfgm}:
\begin{equation}\label{eq:regla_ito}
  df \;=\; \Bigl(\partial_t f + \mu S\, \partial_S f
  + \frac{1}{2}\sigma^2 S^2\, \partial_S^2 f\Bigr)\, dt
  + \sigma S\, \partial_S f\, dW_t ,
\end{equation}
whose only non-classical feature is the term in $\partial_S^2 f$, which
comes from $(dS)^2 = \sigma^2 S^2\,dt$.

\medskip

\noindent\textbf{Under $\Q$.}
Girsanov's theorem \cite[Section~4.4]{tfgm} constructs the risk-neutral
probability of the continuous model: under $\Q$, the process
$\widetilde W_t = W_t + \frac{\mu - r}{\sigma}\,t$ is a Brownian motion,
and
\begin{equation}\label{eq:bs_eds_Q}
  dS_t \;=\; r\, S_t\, dt + \sigma\, S_t\, d\widetilde W_t,
  \qquad
  dX_t \;=\; \Bigl(r - \frac{\sigma^2}{2}\Bigr)\, dt
  + \sigma\, d\widetilde W_t,
\end{equation}
whose second equation is the starting point of
Chapter~\ref{cap:bs_operadores}\enprep. As in \eqref{eq:q_crecimiento_r},
the underlying grows on average at the rate $r$. The drift changes, but
the volatility cannot, for the reason given by the rule
\eqref{eq:regla_dW2}: the sum of the squared increments of $X$ over
$[0, t]$ equals $\sigma^2 t$ with probability one, and two equivalent
probabilities agree on the events of probability one. The volatility
can be read off a single path; the drift cannot be determined from a
path of finite length, and that is why it is the only thing the change
of probability modifies. In Chapter~\ref{cap:bs_operadores}\enprep,
multiplying by an exponential of the log-price removes the drift from
the Hamiltonian: it is a reweighting of paths of the same kind, seen as
a gauge transformation.

\seccionesprevistas{2.5--2.6}
